\section{Common localization and quantum comparisons}
\label{sec:localization}

We next isolate the local variational tools.  A spatial cut records the
electrons removed from a region and leaves an antisymmetric conditional
state for the others.  Its conditional electrostatic field enters an
exact energy identity.  The state may belong to any finite particle-number sector. Because
$E$ is the minimum over all sectors,
we may insert a variable number of trial electrons into the resulting
hole.  Upper and lower semiclassical comparisons will have the same
kinetic coefficient
\[
 \cT=\frac3{10}(3\pi^2)^{2/3}.
\]
An additional kinetic bound, whose coefficient need not be sharp,
controls localization errors.

\subsection{Coulomb energy and a kinetic bound}

Write $\dot H^1(\R^3)$ for the completion of $C_c^\infty(\R^3)$ in
the gradient norm, represented in $L^6(\R^3)$ by Sobolev's inequality.
For a compactly supported real $f\in L^{5/3}(\R^3)$, put
\[
 D(f)=\frac12\int f(K*f).
\]
This definition allows signed $f$.

\begin{lemma}[Coulomb duality and radial smearing]
\label{lem:loc-coulomb}
For compactly supported real $f\in L^{5/3}(\R^3)$, its Coulomb
potential is continuous and belongs to $\dot H^1$, and
\begin{equation}\label{eq:loc-coulomb-duality}
 D(f)=\frac1{8\pi}\|\nabla(K*f)\|_2^2,
 \qquad
 \left|\int fv\right|\leq C\sqrt{D(f)}\|\nabla v\|_2
 \quad(v\in\dot H^1).
\end{equation}
The potential and energy assertions also hold for a nonnegative
$f\in L^{5/3}$ whose distribution belongs to $(\dot H^1)^*$, with
$K*f$ defined by its nonnegative convolution.  Its dual pairing with
any $v\in\dot H^1$ is the absolutely convergent integral $\int fv$.

If $\eta$ is a radial nonnegative probability density supported in
$B(0,b)$, then $K*\eta\leq K$, with equality outside that ball.
In particular, $D(\eta*\rho)\leq D(\rho)$ for every nonnegative
compactly supported $\rho\in L^{5/3}$.
\end{lemma}

\begin{proof}
For compactly supported $f\in L^{5/3}$ we also have $f\in L^{6/5}$.
Riesz representation on $\dot H^1$ gives a unique solution $u$ of
$-\Delta u=4\pi f$.  The convolution $K*f$ is continuous: near the
singularity use $K\in L^{5/2}_{\mathrm{loc}}$ and continuity of
translations in that space; away from it use dominated convergence.
It is bounded locally and decays like $O(|x|^{-1})$ at infinity,
hence belongs to $L^6$.  It has the same distributional Laplacian as
$u$.  A harmonic $L^6$ function is zero, by mollification and the
mean-value property on arbitrarily large balls.  Thus $u=K*f$.
Testing its weak equation with $u$ gives the first identity in
\eqref{eq:loc-coulomb-duality}; testing with $v$ and using
Cauchy--Schwarz gives the second.

For the noncompact nonnegative case, positivity first gives
\[
 \int f|\varphi|\leq\|f\|_{(\dot H^1)^*}\|\nabla\varphi\|_2
 \qquad(\varphi\in C_c^\infty).
\]
To obtain this estimate, approximate the absolute value by smooth
convex truncations with derivative bounded by one.  Completion and
an almost-everywhere convergent subsequence extend the estimate to
$\dot H^1$ and identify its dual pairing with actual integration.
Let $u$ be the Riesz solution, and let
$u_j=K*(\mathbf1_{B(0,j)}f)$.  Testing the difference equation with
$(u_j-u)_+$ gives $u_j\leq u$.  These positive-part tests belong to
$\dot H^1$: an $L^6$ function with gradient in $L^2$ is in the
completion by large-radius cutoff and mollification.  For the cutoff
error, H\"older's inequality bounds the gradient of the cutoff times
the function by its $L^6$ norm on the escaping annulus.  Now
$u_j\uparrow K*f$ and $u_j\leq u\in L^6$, so dominated convergence
in $L^6$, the distributional equation, and harmonic uniqueness give
$K*f=u$.  Its energy identity follows by testing with $u$.
The convolution is continuous also in this case.  Choose a point $x_0$
where its value is finite.  On any fixed bounded set of evaluation points,
the sufficiently distant kernel tails are dominated by a constant times
$K(x_0-\cdot)f$, which is integrable.  Dominated convergence handles that
tail, and the local $L^{5/3}$--$L^{5/2}$ argument handles the remaining
compact part.

The spherical average of $K$ at distance $d$ from a sphere of radius
$s$ is $1/\max(d,s)$, by direct angular integration.  Decomposing a
radial probability into spherical averages proves the smearing
assertion.  The difference of two independent radial displacements
is again radial; applying the same assertion to its distribution
and integrating against $\rho(x)\rho(y)$ proves the energy inequality.
\end{proof}

For a positive trace-class operator $\gamma$ on
$L^2(\R^3;\mathbb C^2)$, write $\rho_\gamma$ for its spatial density,
with spins summed.  Its kinetic energy is
$\operatorname{Tr}(-\Delta/2)\gamma$, defined by the nonnegative form.
We use the Lieb--Thirring kinetic inequality \citep{LiebThirring1975}
in the following non-sharp form.  The short proof by spectral projections
also follows the distribution-energy approach of \citet{Rumin2011}.

\begin{lemma}[Occupation and kinetic energy]
\label{lem:loc-kinetic}
If $0\leq\gamma\leq1$ has finite trace, then
\begin{equation}\label{eq:loc-coarse-kinetic}
 \operatorname{Tr}(-\Delta/2)\gamma
 \geq c\int\rho_\gamma^{5/3}
\end{equation}
for a universal $c>0$.  The one-body density matrix of an
antisymmetric $k$-particle vector $f$ is bounded above by
$\|f\|_2^2$ times the identity.  The same bound holds after
integration in auxiliary variables, and such matrices may be summed
whenever the corresponding squared norms sum to one.
\end{lemma}

\begin{proof}
For a unit spin orbital $v$, the number operator
$\sum_{i=1}^k|v\rangle\langle v|_i$ is a projection on the
antisymmetric space: in a Slater basis containing $v$, its eigenvalue
is zero or one.  Its expectation is at most $\|f\|_2^2$.  This
proves the occupation assertion and its versions with auxiliary
integration and mixtures.

For the kinetic bound, take an eigen-expansion
$\gamma=\sum_\alpha\lambda_\alpha
 |u_\alpha\rangle\langle u_\alpha|$, and let $P_s$ be Fourier
projection onto $|p|^2/2\leq s$.  The density of its low-frequency
parts is at most $Cs^{3/2}$: Bessel's inequality bounds each spin
evaluation against the corresponding Fourier vector, and there are
two spins.  The triangle inequality in the weighted orbital-index
and spin norm gives, almost everywhere,
\[
 \sum_{\alpha,\sigma}\lambda_\alpha
       |((1-P_s)u_\alpha)(x,\sigma)|^2
 \geq\big(\sqrt{\rho_\gamma(x)}-Cs^{3/4}\big)_+^2.
\]
Integrating the left side in $x$ and then in $s>0$ gives the kinetic
trace by Plancherel and the layer-cake formula.  The integral of the
right side in $s$ is a positive constant times
$\rho_\gamma(x)^{5/3}$.  Truncation proves the assertion also when
the kinetic trace is infinite.
\end{proof}


\begin{lemma}[Global atomic density at an energy offset]
\label{lem:atomic-global-density}
For one nucleus of charge $Z\ge1$, every $\delta$-state in a sector
$n\le3Z$ satisfies
\begin{equation}\label{eq:atomic-global-density}
 \int\rho_\psi^{5/3}\le C(Z^{7/3}+\delta).
\end{equation}
\end{lemma}
\begin{proof}
Drop the nonnegative pair repulsion, use $E\le0$, and split attraction
at $R=Z^{-1/3}$. Outside that ball it is at most $Zn/R\le3Z^{7/3}$.
Inside, Young's inequality gives
\[
 \int_{|x|<R}\frac Z{|x|}\rho_\psi
 \le\frac c2\int\rho_\psi^{5/3}
       +C\int_{|x|<R}(Z/|x|)^{5/2}
 \le\frac c2\int\rho_\psi^{5/3}+CZ^{7/3},
\]
where $c$ is the constant in Lemma~\ref{lem:loc-kinetic}.
Combine this with that kinetic lower bound and the energy bound
$q_n[\psi]\le E+\delta\le\delta$.
\end{proof}

\subsection{Cuts, conditional states, and the exact energy identity}

The underlying tools are the IMS formula
\citep[Lemma~2.4]{Solovej2003} and localization into correlated
particle-number sectors \citep[Proposition~3.1 and Example~3.3]{Lewin2011Geometric}.
Here we retain the removed coordinates as data, so that their conditional
core field remains available in the exact energy identity.

Consider a finite sequence of deterministic nonnegative Lipschitz
functions $(s_h,c_h)$ on $\R^3$, with $s_h^2+c_h^2=1$.  At stage $h$
the factor $s_h$ removes particles from what remains of the core.
Define the first-out factors and the final core factor by
\begin{equation}\label{eq:loc-first-out}
 p_h=s_h\prod_{l<h}c_l,\qquad c=\prod_hc_h,
 \qquad o=\Big(\sum_hp_h^2\Big)^{1/2}.
\end{equation}
Thus $o^2+c^2=1$.  For a normalized $\psi\in\mathcal Q_n$ and
$I\subset\{1,\ldots,n\}$, put
\[
 \psi_I=\prod_{i\in I}o(x_i)\prod_{i\notin I}c(x_i)\psi.
\]
Under the squared laws of these functions, whose total mass is one,
record $I$ and the positions and spins of its out particles.  Decorate
each such particle with stage $h$ with probabilities $p_h^2/o^2$ at
its position.  Denote the resulting data by $X$.

Given $X$, normalize the slice in the remaining core variables.
It is an antisymmetric form-domain state almost surely.  Write $e_X$
for its energy, $\rho_X^c$ for its one-body density, and
\[
 \Phi_X=V-K*\rho_X^c.
\]
For the vacuum slice set $e_X=0$ and $\rho_X^c=0$.

\begin{lemma}[Conditional localization identity]
\label{lem:loc-identity}
With the preceding definitions, set
\[
 T_o=\frac12\sum_I\sum_{i\in I}\|\nabla_i\psi_I\|_2^2.
\]
There is a nonnegative localization error satisfying
\begin{equation}\label{eq:loc-error}
 \mathrm{IMS}\leq\frac12\E_\psi\sum_i\sum_h
       \big(|\nabla s_h|^2+|\nabla c_h|^2\big)(x_i)
\end{equation}
such that
\begin{equation}\label{eq:loc-cut}
 q_n[\psi]+\mathrm{IMS}
 =\E e_X+T_o-\E\sum_{i\ \mathrm{out}}\Phi_X(x_i)
                   +\E\sum_{i<j\ \mathrm{out}}K(x_i-x_j).
\end{equation}
The out kinetic energy satisfies
\[
 T_o\geq c\int(o^2\rho_\psi)^{5/3},
\]
where $\rho_\psi$ is the raw one-body density.  Conditional quantities
have versions jointly measurable in the data and the field evaluation
point.  Appending deterministic cuts admits a coupling in which these
decorated data are nested, and hence conditional expectations of raw
quantities satisfy the tower property.
\end{lemma}

\begin{proof}
Successive squared-gradient splitting gives
\[
 \sum_h|\nabla p_h|^2+|\nabla c|^2
 =\sum_h\Big(\prod_{l<h}c_l^2\Big)
                 (|\nabla s_h|^2+|\nabla c_h|^2).
\]
The norm-map inequality bounds $|\nabla o|^2$ by
$\sum_h|\nabla p_h|^2$.  Expanding the product partition in the
kinetic form, the mixed derivatives cancel because
$o\nabla o+c\nabla c=0$.  This proves the stated IMS bound.
Multiplication potentials split exactly by the squared partition.

Fix the recorded subset $I$, and write $u$ for its out positions and
spins and $\xi$ for the remaining positions and spins.  Put
\[
 \kappa_I(u)=\int|\psi_I(u,\xi)|^2\,\mathrm d\xi,
 \qquad
 \chi_X(\xi)=\frac{\psi_I(u,\xi)}{\sqrt{\kappa_I(u)}}
 \quad\text{when }\kappa_I(u)>0,
\]
where the integral includes the core spin sum.  Thus $\chi_X$ is the
normalized core vector, with its phase retained.  The stage decoration
depends only on $u$ and does not change this normalized slice.
Fubini gives its $H^1$ regularity almost surely; permutations of core
variables preserve the multiplying factors and therefore its antisymmetry.
The conditional core--out repulsion is exactly
\[
 \sum_{i\ \mathrm{out}}\int K(x_i-x)\rho_X^c(x)\,\mathrm dx.
\]
Grouping the core form, out kinetic form, and the remaining
multiplication potentials proves \eqref{eq:loc-cut}.  All absolute
Coulomb terms and kinetic terms are integrable in this calculation
by the form bounds and the finite partition.  The slice ratios and
their integrals give the stated measurable versions; assign arbitrary
values on null slices.

For each $I$, apply the occupation assertion of
Lemma~\ref{lem:loc-kinetic} to the out group, integrating the other
variables.  Its one-body matrix $\gamma_I^o$ obeys
$\gamma_I^o\leq\|\psi_I\|_2^2$.  Thus
$\gamma_o=\sum_I\gamma_I^o\leq1$, even though the out particle number
varies with $I$.  Its kinetic trace is $T_o$, and its density is
$o^2\rho_\psi$, by summing the squared partition with one label
required to be out.  Equation~\eqref{eq:loc-coarse-kinetic} proves the
kinetic assertion.

Finally, the entire construction has a single sampling description.
Conditional on the raw configuration, independently label each
coordinate by its first-out stage or by final core, with probabilities
$p_h^2$ and $c^2$.  An appended cut only splits the previous core
probability.  Forgetting its new-stage labels and their variables
recovers the old observation.  This is the required nesting and
proves the tower assertion.
\end{proof}

\subsection{Insertion of trial electrons}

We use the coherent-packet construction
\citep[Lemma~8.2]{Solovej2003}. The next lemma turns a nonnegative
density into a genuine fermionic trial state.  Its particle number is random, with integer values;
the global sector minimum in \eqref{eq:sector-global-minimum} is what
makes this permissible.

Fix once and for all a real smooth radial $g$ supported in the unit
ball, with $\int g^2=1$, and set $g_b(x)=b^{-3/2}g(x/b)$.

\Needspace{4\baselineskip}
\begin{lemma}[Trial particles in a hole]
\label{lem:loc-hole-trial}
Fix a conditional core slice with finite energy $e_X$ and field
$\Phi_X$ as above.  Let $\rho\geq0$ belong to $L^{5/3}$ and have
compact support.  Suppose its packet region
$\operatorname{supp}\rho+\overline{B(0,b)}$ has positive distance
from the spatial support allowed to the core particles.  Then
\begin{equation}\label{eq:loc-trial}
 E\leq e_X+\cT\int\rho^{5/3}
       +Cb^{-2}\int\rho
       -\int\Phi_X(g_b^2*\rho)+D(\rho).
\end{equation}
The constant depends only on the fixed function $g$.  The vacuum
trial also gives $E\leq e_X$.
\end{lemma}

\begin{proof}
For each center $y$, momentum $p$, and spin $\sigma$, use the orbital
\[
 v_{y,p,\sigma}(x,\tau)
       =g_b(x-y)e^{ip\cdot x}\mathbf1_{\{\tau=\sigma\}}.
\]
Integrate its projector over both spins and
$|p|\leq(3\pi^2\rho(y))^{1/3}$ with measure
$\mathrm dy\,\mathrm dp/(2\pi)^3$, obtaining an operator $\gamma$.
The unrestricted packet resolution is the identity: apply
Plancherel in $p$ and then integrate $g_b^2$ in $y$.  Restriction
therefore gives $0\leq\gamma\leq1$.  Momentum-ball integration gives
\begin{align}
 \operatorname{Tr}\gamma&=\int\rho,
 &\rho_\gamma&=g_b^2*\rho,\notag\\
 \operatorname{Tr}(-\Delta/2)\gamma
   &=\cT\int\rho^{5/3}
        +\frac12\|\nabla g\|_2^2b^{-2}\int\rho.
 \label{eq:loc-packet-energy}
\end{align}
The reality of $g$ eliminates the kinetic cross term.  These traces
are finite.  Every positive-eigenvalue orbital of $\gamma$ belongs
to $H^1$ and is supported in the packet region.

Take a finite eigen-truncation of $\gamma$.  Occupy each retained
orbital independently with Bernoulli probability equal to its
eigenvalue.  Each realization is a normalized Slater determinant,
including the possible vacuum.  The average one-body energy is the
trace energy.  Its average repulsion is the direct energy of the
truncated density minus
\[
 \frac12\sum_{\sigma,\tau}\iint
       K(x-y)|\gamma(x,\sigma;y,\tau)|^2\,\mathrm dx\,\mathrm dy,
\]
where here $\gamma$ denotes the finite truncation.  In particular
the average repulsion is at most the direct energy.  Equal-orbital
direct and exchange contributions cancel, as required by the
Bernoulli occupation rule.

Wedge each realization with the normalized core slice.  Because
their spatial supports are separated, the terms assigning the two
groups to different sets of particle labels have disjoint supports.
Normalized antisymmetrization consequently preserves the sum of the
group energies and adds exactly their direct cross repulsion.  Each
resulting vector belongs to its own antisymmetric particle-number
sector and has energy at least $E$.  Averaging gives the desired
inequality with the finite-truncation density and kinetic energy.

For completeness, passage to the full operator needs no selection
of infinite Slater states.  The truncated densities increase almost
everywhere to $g_b^2*\rho$, and their kinetic energies increase to
\eqref{eq:loc-packet-energy}.  The full density is bounded and
compactly supported, so nuclear integrals converge.  The core
potential is bounded on the packet region by support separation,
so its integrals converge as well.  Direct energies converge
monotonically.  Lemma~\ref{lem:loc-coulomb} gives
$D(g_b^2*\rho)\leq D(\rho)$, proving \eqref{eq:loc-trial}.
The same argument with no inserted orbitals proves $E\leq e_X$.
\end{proof}

The inequality is pointwise for every admissible core slice outside
a null set.  Later we choose $\rho$ as a measurable function of its
conditional field.  Only the displayed functional is then averaged;
no measurable choice of Slater determinants is necessary.

\subsection{A lower comparison with the same kinetic coefficient}

An arbitrary out configuration also determines a bounded density,
obtained by spreading each retained particle over a small radial
kernel.  To retain the coefficient $\cT$ in the lower bound, use
Neumann eigenvalues in cubes and then average the grid. This is the
Neumann-bracketing approach to the semiclassical kinetic coefficient
\citep[Theorems~III.10--III.13]{LiebSimon1977}; the retained subset and
its radial fine density are specified below.

For $b>0$, define the radial probability kernel
\begin{equation}\label{eq:loc-fine-kernel}
 \vartheta_b(x)=\int_{\mathrm{SO}(3)}
       b^{-6}\prod_{j=1}^3\big(b-|(Qx)_j|\big)_+\,\mathrm dQ,
\end{equation}
where Haar measure has mass one.  It is bounded by $Cb^{-3}$ and
supported in $\overline{B(0,\sqrt3b)}$.

\begin{lemma}[Fine density and lower energy bounds]
\label{lem:loc-cubes}
In the setting of Lemma~\ref{lem:loc-identity}, let $n_o$ be the
out particle count.  Choose, measurably, any subset of the out
particles and set
\[
 \sigma(x)=\sum_{i\ \mathrm{retained}}\vartheta_b(x-x_i).
\]
Then
\begin{align}
 T_o&\geq\E\left[\cT\int\sigma^{5/3}
               -Cb^{-2}(n_o^{4/3}+n_o)\right],
 \label{eq:loc-fine-kinetic}\\
 \sum_{i<j\ \mathrm{out}}K(x_i-x_j)
   &\geq D(\sigma)-Cn_o/b
 \quad\text{almost surely}.
 \label{eq:loc-fine-pair}
\end{align}
\end{lemma}

\begin{proof}
On a cube of side $b$, the spin-two Neumann eigenvalues are
$\pi^2|m|^2/(2b^2)$, for $m\in\mathbb N_0^3$, each with two spin
states.  The number of lattice points of radius at most $s$, with
this multiplicity, is bounded by
\[
 \frac\pi3(s+\sqrt3)^3:
\]
attach a disjoint unit cube to each octant lattice point and enclose
their union in the octant ball of radius $s+\sqrt3$.  Thus the $j$th
ordered radius is at least
$((3j/\pi)^{1/3}-\sqrt3)_+$.  Summing their squared radii and using
$\sum_{j\leq n}j^{2/3}\geq(3/5)n^{5/3}$ gives the lower bound
\begin{equation}\label{eq:loc-neumann-levels}
 b^{-2}\big(\cT n^{5/3}-C(n^{4/3}+n)\big)
\end{equation}
for the sum of the first $n$ eigenvalues.  The leading coefficient
is exactly $\cT$.

Partition the out coordinates into assignments to the cubes of a
grid.  Restriction to a cube is admissible for its Neumann form;
no zero extension across its boundary is involved.  Within each
cube the corresponding variables remain antisymmetric.  The
occupation bound, now in the Neumann basis, gives
\eqref{eq:loc-neumann-levels} after slicing all auxiliary variables.
If $n_C$ denotes the out count in a cube, the leading term is
\[
 \cT\int\left(\sum_C\frac{n_C}{b^3}\mathbf1_C(x)\right)^{5/3}
                  \mathrm dx.
\]
The error is at most $Cb^{-2}(n_o^{4/3}+n_o)$, since
$\sum_Cn_C^{4/3}\leq n_o^{4/3}$.

Average the grid over a uniform translation in a fundamental cube
and over rotations.  At a fixed out configuration, its averaged
cell density is $\sum_{i\ \mathrm{out}}\vartheta_b(x-x_i)$.
Before rotation, the probability that two points are in the same
translated cube gives precisely the triangular product in
\eqref{eq:loc-fine-kernel}.  Jensen's inequality bounds the power
integral from below by the power of this averaged density.
Dropping any chosen particles can only decrease that density and
its nonnegative power.  This proves \eqref{eq:loc-fine-kinetic}.

For the pair bound, first discard all pairs not both retained.
Smearing each retained point by $\vartheta_b$ decreases every
distinct pair interaction by Lemma~\ref{lem:loc-coulomb}, because
the difference of two independent radial displacements is radial.
Each self energy is bounded by $C/b$, using the kernel's support and
$L^\infty$ bound.  Subtracting these self energies from $D(\sigma)$
proves \eqref{eq:loc-fine-pair}.
\end{proof}

We now have both sides of the local comparison.  The upper trial
uses an arbitrary nonnegative density in a hole; the lower bound
uses the fine density of the electrons actually removed.  Their
matching $\cT\int\rho^{5/3}$ terms will identify the same convex
Thomas--Fermi functional, while Lemma~\ref{lem:loc-coulomb} measures
the difference of the two densities.
